
%\needspace{4\baselineskip}
\color{uniblue}{\subsection{Восстановление нормального режима (ВНР)}}
\color{black}

\begin{enumerate}

\item
ВНР предназначено для автоматического восстановления схемы доаварийного режима после действия АВР. Устройство может выполнять обратное переключение, что подразумевает отключение секционного выключателя (выключателя резервного ввода) и включение основного рабочего ввода. В устройстве восстановление нормального режима представлено функцией «ВНР».

\item
В алгоритме функции предусмотрено два режима работы: без перерыва питания потребителей (ВВ-СВ) и с перерывом питания (СВ-ВВ). В первом случае работа осуществляется без перерыва питания, включается основной рабочий ввод и потом отключается секционный выключатель (выключатель резервного ввода). Во  втором  случае  сначала  отключается  секционный  выключатель,  затем  проверяется  его  включенное состояние и, после этого включается выключатель основного рабочего ввода. Данные режимы выбираются уставкой SGF2 «Очередность\_ВНР» (уставка программного переключателя  SGF2  по  умолчанию выставлена в положение «СВ-ВВ»). 

\item
Критерием пуска алгоритма ВНР служит появление рабочего напряжения на ТН до ввода. Алгоритм ВНР может работать только после срабатывания АВР. Действие функции ВНР блокируется, если выведена (заблокирована) функция АВР. 

\item
Готовность  схемы  ВНР  и  активация  соответствующего  сигнала «ВНР:  Ожидание  ВНР» собирается после активации команды на включение СВ от схемы АВР и включенного положения секционного выключателя с дополнительным контролем отсутствия напряжения на какой-либо секции. При выполнении всех перечисленных условий формируется выходной сигнал ожидания  ВНР, который сохраняется  в энергонезависимой памяти до появления сигналов для сброса триггера.

\item
Описание алгоритмов работы ВНР:
\begin{enumerate}[label=\arabic*), noitemsep]
    \item Алгоритм работы схемы ВВ-СВ (ВНР без перерыва питания):
    \begin{itemize}
        \item При действии ВНР без перерыва питания, необходимо выставить уставку SGF2 в положение «ВВ-СВ» (т.е. Вводной выключатель \rightarrow  Секционный выключатель). В этом режиме сначала формируется команда на включение выключателя ввода, АВР которой выполнялся, а затем выполняется команда на отключение СВ. Время ввода ВНР задается таймером Т1. 
        \item Для включения вводного выключателя первой секции шин необходимо запараметрировать на выходное реле сигнал «ВНР: Включение ВВ 1сш». Аналогично выполнить для включения вводного выключателя для второй секции шин.  
        \item Для начала работы ВНР необходимо, чтобы было восстановлено питающее напряжение до вводного выключателя, то есть появился сигнал наличия напряжения «КНН до ввода 1сш» или «КНН до ввода 2сш». Отсутствие сигнала АВР. Вводной выключатель первой или второй секции шин отключен («ВВ 1сш отключен» или «ВВ 2сш  отключен»). При  наличии  готовности  к  включению  вводного выключателя первой или второй секции шин, по истечении времени таймера Т2, будет сформирована команда на включение соответствующего вводного выключателя.
        
    \end{itemize}
    \item Алгоритм работы схемы СВ-ВВ (ВНР с перерывом питания):
    \begin{itemize}
        \item При действии ВНР с перерывом питания, необходимо выставить уставку SGF2 в положение «СВ-ВВ» (т.е. Секционный выключатель \rightarrow Вводной выключатель). В этом режиме сначала формируется команда на отключение секционного выключателя путем воздействуя на себя, а затем формируется команда  на включение  выключателя  той  секции,  АВР  которой  была  осуществлена. 
        \item Для работы в данном режиме необходимы те же условия логики, что описаны выше. При наличии готовности к включению вводного выключателя, по истечению выдержки времени таймера Т2 формируется команда на отключение СВ «ВНР: Отключить». После отключения выключателя в логику ВНР поступает сигнал  отключенного  положения  СВ  «В:  отключен»  и  по  истечению  выдержки  времени  таймера  Т3 формируется  команда  на  включение  ВВ  данной  секции  и  сброс  триггеров,  если  сохранились  условия готовности к включению выключателя.
    \end{itemize}
\end{enumerate}

\item
Ввод функции в работу на этапе параметрирования устройства осуществляется программным переключателем SGF1 «Ввод\_функции», при этом необходимо введенное состояние функции АВР (активный сигнал «АВР: Ввод»). Информация о введенной функции отображается выходным сигналом «ВНР: Ввод».

\item\sloppy
Функция может быть выведена из работы оперативно путем активации сигнала «ОВ ВНР». Выведенное состояние функции характеризуется активным состоянием сигнала «ВНР: Оперативный вывод».

\item
Алгоритм работы функции выполнен в соответствии с рисунком \ref{vnr:img1}. В таблице \ref{vnr:tbl1} приведены параметры, необходимые для настройки функции.

\vspace{3mm}
\begin{figure}[H]
\centering
\includegraphics[width=1\textwidth,height=1\textheight,keepaspectratio]{\fbpath/411. ВНР СВ/_latex/img/vnr.pdf}
\captionof{figure}{Функциональная схема логики ВНР}\label{vnr:img1}
\end{figure}

\end{enumerate}
